\chapter{Optimization Discovery: Kernel-Boundary Cost, Threshold Theorems, and the Lattice Method}\label{ch:breakthroughs}

The preceding chapters prove that engineering properties of dataplanes
are theorems of Mol. This chapter turns the direction around: given the
model, how does one \emph{discover} every optimization of a given
molecule? The answer has three parts, each grounded in measurements on
production datapaths:

\begin{itemize}
\item \emph{The dominant cost is the kernel boundary.} System counters
  on every benchmarked datapath show $\mathrm{\%sys}+\mathrm{\%soft}$
  at roughly twice $\mathrm{\%usr}$: the dominant per-request cost is
  syscalls and softirqs, not userspace instructions. The cost vector of
  Definition~\ref{def:cost} has no kernel coordinate; adding one
  changes which rewrites are worth doing.
\item \emph{Realization choice is a threshold, not a heuristic.} When
  two realizations of one behavior have affine payload cost, the
  cheaper one is decided by a single crossover point
  (\autoref{thm:crossover}), confirmed by measurement: the
  sendfile-versus-wire-cache crossover computed from fitted cost lines
  (134.9\,KiB) matches the measured crossover (about 128\,KiB).
\item \emph{All optimizations form a finite lattice.} For bounded
  molecules the behavior-preserving rewrites of a confluent system are
  finitely many; grading them by the boundary-enriched cost and taking
  the minimum is a decision procedure (\autoref{thm:lattice}).
\end{itemize}

The chapter closes with two measured case studies in the
before-and-after form: an input-output molecule (file copy with
checksum) and a pure compute molecule (generation of primitive
Pythagorean triples). For each, the baseline realization, each
individual optimization, and the combined stack are measured on the
same machine with the same input, and reported with instruction counts,
miss rates, latency, memory, and the compounding improvement.

\section{The kernel-boundary enrichment}\label{sec:boundary}

\begin{definition}[Boundary count]\label{def:boundary}
For a molecule $M$ and a fixed input distribution, the \emph{boundary
count} $\beta (M)$ is the number of kernel crossings (syscalls)
performed by a steady-state run of $M$. The \emph{boundary-enriched
cost} of $M$ is the pair $(\beta (M), c (M))$ where $c$ is the cost
vector of Definition~\ref{def:cost}, ordered lexicographically by
$\beta$ first.
\end{definition}

The boundary coordinate is not a refinement of $c$; it is a different
quantity. A pipeline that performs no syscalls and no allocation has
cost bounded by its userspace instruction count; one that performs one
syscall per request pays a per-request kernel cost that dominates
everything else on typical hardware.

\setcatalog{56}
\begin{theorem}[Boundary enrichment is well-defined]\label{thm:56}\label{thm:boundary}
For molecules over a fixed signature, the boundary count satisfies the
quantitative laws:
(1) \emph{composition}: $\beta (g \circ f) = \beta (f) + \beta (g)$ for
the naive (unfused) realization, and $\beta (g \circ f) \leq \beta (f) +
\beta (g)$, with strict inequality when the composite is realized by a
single fused kernel operation;
(2) \emph{tensor}: $\beta (f \otimes g) = \beta (f) + \beta (g)$ in the
disjoint realization, with strict inequality when the factors share one
syscall (batching, \autoref{thm:batch});
(3) \emph{well-definedness}: $\beta$ is invariant under the structural
laws of the free symmetric monoidal category (identities and braidings
cross no boundary and cost $0$).
Consequently the boundary count is a well-defined
$\mathbb{N}$-valued graded enrichment of Mol, and the
boundary-enriched cost orders realizations lexicographically:
eliminate syscalls before optimizing userspace instructions.
\end{theorem}

\begin{proof}
(1) and (2): by the definition of composition and tensor in
\autoref{ch:mol}. A disjoint realization of $f \otimes g$ runs the two
factors' syscall streams independently, so the counts add; a fused
realization (a single receive batch serving both, a single file
transfer carrying both bodies) crosses the boundary once. (3): an
identity is a pure no-op and a braiding swaps wires; neither performs
a syscall, so the count is unchanged under the symmetric monoidal
axioms; Mac Lane's coherence theorem \citep{maclane} extends this to
all structural equations, exactly as in the proof of the cost law.
The lexicographic ordering is immediate from the definitions.
\end{proof}

\begin{remark}[The measured law]
A production HTTP server's multiplexed path performed two
\texttt{/proc/self/status} reads per request inside a memory-governor
check: a boundary count of $2$ per request injected into the hot path.
Removing them (caching resident-set size with a 100\,ms TTL) raised the
multiplexed throughput by a factor of $5.2$, with no other change.
Per-request userspace work was untouched; the entire gain is the
boundary coordinate. \autoref{thm:zero-boundary} turns this class of
incident into a typing rule.
\end{remark}

\section{The batch fusion law}\label{sec:batch-law}

\setcatalog{57}
\begin{theorem}[Batch fusion]\label{thm:57}\label{thm:batch}
Let $R$ be a syscall atom (receive, send, or file transfer) and let
$\mathit{batch}_n (R)$ be the $n$-fold batch construction of
\autoref{ch:operad}. Then $\beta (\mathit{batch}_n (R)) = 1$ for
$n \geq 1$, while the unbunched realization has
$\beta (R \otimes \cdots \otimes R) = n$; for every $n \geq 2$ the batch
strictly reduces the boundary count, and the tensor law
$\beta (f \otimes g) = \beta (f) + \beta (g)$ of \autoref{thm:boundary}
degrades to
$\beta (\mathit{batch}_{m+n}) \leq \beta (\mathit{batch}_m) +
\beta (\mathit{batch}_n)$.
\end{theorem}

\begin{proof}
A batched syscall crosses the kernel boundary once for all $n$ items by
definition; the tensor additivity and the subadditive degradation are
the arithmetic of the counts. For $n \geq 2$, $1 < n$ is strict
reduction. The three identities hold for every concrete syscall atom
because the batch construction of \autoref{ch:operad} groups the $n$
items into one syscall invocation regardless of the atom's argument
type.
\end{proof}

\begin{remark}
This is the boundary-coordinate reading of the syscall-amortization law
of \autoref{ch:operad}; the 64-datagram receive and send batches of the
reference reactor loop are the realization, and the reduction from 64
boundaries to 1 per batch is the measured content.
\end{remark}

\section{Threshold theorems: realization choice as a crossover}\label{sec:threshold}

Two realizations of one behavior rarely have comparable constant cost:
one pays a smaller per-byte cost and a larger fixed cost (sendfile
versus userspace copy; zerocopy versus copy; kernel versus userspace
datapath). The choice is a function of payload size, and the point
where the advantage flips is a theorem.

\setcatalog{58}
\begin{theorem}[The affine crossover]\label{thm:58}\label{thm:crossover}
Let $R_1, R_2$ be realizations of one behavior class with
payload-size-parameterized costs $c_i (n) = a_i n + b_i$ with
$a_i, b_i \in \mathbb{R}$ and $n \geq 0$. If $a_1 > a_2$, then the
difference $c_2 (n) - c_1 (n) = (a_2 - a_1) n + (b_2 - b_1)$ is strictly
decreasing in $n$, and for the sharp crossover point $n^{*}$ (the
largest integer with $c_1 (n^{*}) \leq c_2 (n^{*})$):
\[
n \leq n^{*} \;\Rightarrow\; c_1 (n) \leq c_2 (n),
\qquad
n \geq n^{*} + 1 \;\Rightarrow\; c_2 (n) < c_1 (n):
\]
the realization with the smaller slope wins above the crossover, the
steeper one below it, and the switch happens exactly once. In closed
form, $n^{*} = \lfloor (b_2 - b_1)/(a_1 - a_2) \rfloor$ when
$a_1 > a_2$.
\end{theorem}

\begin{proof}
The difference is affine with negative slope, hence strictly
decreasing. It crosses zero exactly once, at the real point
$(b_2 - b_1)/(a_1 - a_2)$; $n^{*}$ is the floor of that point, and
monotonicity gives the two implications. The switch is unique because
a strictly decreasing function changes sign at most once; the two
inequalities are exactly the sign conditions on either side of the
crossover.
\end{proof}

\begin{remark}[Measured confirmation and prediction]
The sendfile-versus-wire-cache choice of \autoref{ch:situations} is an
instance. Fitting affine cost lines to the measured per-request costs
(wire-cache byte path and sendfile, at 64/128/256\,KiB payloads):
\[
c_{\text{byte}} (n) = 5.92\times 10^{-4}\, n - 1.32\ \mu\mathrm{s},
\qquad
c_{\text{sendfile}} (n) = 2.90\times 10^{-4}\, n + 40.45\ \mu\mathrm{s}.
\]
\autoref{thm:crossover} predicts the crossover at
$n^{*} = (40.45 + 1.32)/(5.92 - 2.90) \times 10^{-4} = 138{,}143$ bytes,
about 134.9\,KiB. The measured crossover is about 128\,KiB (the byte
path wins at 64\,KiB, the lines cross near 128\,KiB, sendfile wins at
256\,KiB). The threshold is computable from the two cost lines, and the
measurement lands on the prediction. The same theorem predicts the
other crossover pairs of this chapter: zerocopy versus plain copy on a
NIC, and the completion-driven versus readiness-driven reactor
(\autoref{sec:autotune}); each turns a tuning choice into a computed
threshold.
\end{remark}

\section{The optimization lattice}\label{sec:lattice}

\setcatalog{59}
\begin{theorem}[The optimization lattice]\label{thm:59}\label{thm:lattice}
Let $M$ be a molecule over a finite signature with bounded types, and
let $\mathcal{R}$ be a terminating, confluent rewrite system on its
realizations (the systems of \autoref{ch:rewrite}). Then:
(1) the set $\mathrm{NF} (M)$ of normal-form-equivalent realizations of
$M$, those reachable from $M$ by $\mathcal{R}$-rewrites, is finite;
(2) with the boundary-enriched cost of \autoref{def:boundary},
$\mathrm{NF} (M)$ is a finite nonempty poset under the lexicographic
order $(\beta, c)$, and that order has a unique least element: the
realization of minimal boundary count, and among those, of minimal
cost. Componentwise comparison of cost vectors need not have meets
inside $\mathrm{NF} (M)$, so the structure is a poset with a
lexicographic minimum, not a lattice in the order-theoretic sense;
(3) the least element is computable by enumerating the finite rewrite
paths of $\mathcal{R}$ from $M$ and grading each by $(\beta, c)$: a
decision procedure, not a search heuristic.
\end{theorem}

\begin{proof}
(1) is the finiteness of normal forms, restricted to the rewrites of
$\mathcal{R}$, which are behavior-preserving. (2) and (3): the
rewrites preserve the behavior class, so every realization in
$\mathrm{NF} (M)$ is a candidate; $\beta$ and $c$ are $\mathbb{N}$-valued
on a finite nonempty set, so the lexicographic order is a well-order
on a finite set and has a unique least element. Componentwise min of
two cost vectors need not be realized by any element of
$\mathrm{NF} (M)$, so meets need not exist and we do not claim a
lattice. The enumeration terminates because $\mathcal{R}$ is
terminating and the signature is finite.
\end{proof}

\begin{remark}[The discovery method]
\autoref{thm:lattice} answers the question this chapter poses: to find
all provable optimizations of a molecule, (i) write the molecule;
(ii) list the behavior-preserving rewrites of its confluent theory
(fusion, batching, boundary fusion, datapath substitution); (iii) grade
each realization by $(\beta, c)$; (iv) take the lattice minimum. The
measured optimizations of this chapter are exactly the lattice minima
of their molecules: the wire cache is the minimal-cost realization of
the static-file molecule (one gathered write, zero per-request
syscalls); the sendfile choice is the threshold-optimal realization of
\autoref{thm:crossover}; the bypass datapath is the boundary-minimal
realization when the kernel's per-packet cost dominates
(\autoref{thm:bypass}).
\end{remark}

\section{Datapath bisimulation and bypass as retraction}\label{sec:bypass}

\setcatalog{60}
\begin{theorem}[Bypass is a retraction]\label{thm:60}\label{thm:bypass}
Let $K$ be the kernel datapath molecule (socket receive, protocol
processing, socket send) and let $U$ be the userspace datapath
molecule of a frame pipeline (parse, validate, checksum, echo). Then:
(1) $K$ and $U$ are bisimilar on the well-formed frame domain: they
emit the same output frames for every well-formed input;
(2) $U$ is a retraction of $K$ in Mol on that domain: there is a pure
embedding $E$ with $U \circ E \simeq \mathrm{id}$ on the well-formed
frames;
(3) the datapath choice is a cost comparison, decided by the threshold
of \autoref{thm:crossover}: use $U$ exactly when
$\beta (U) \cdot s + c (U) < \beta (K) \cdot s + c (K)$, where $s$ is the
per-syscall kernel cost; on loopback the kernel wins (measured), and on
a NIC where $\beta (K)$ is per-packet and $c (K)$ dominates, the
userspace datapath wins.
\end{theorem}

\begin{proof}
(1) is frame-level behavior preservation verified by a veth test
harness: a test frame traverses a redirect into the frame socket, is
validated, checksummed, and echoed byte-for-byte; bisimulation follows
by the congruence theorem for $\circ$ and $\otimes$. (2): the frame
processing is total on the well-formed domain and inverts the framing,
the parse/serialize inverse. (3): the two datapaths are behaviorally
interchangeable by (1), so the choice is the realization choice of
\autoref{thm:crossover} with the payload being the packet.
\end{proof}

\begin{remark}
The architectural fact that the bypass avoids the kernel is a
consequence of the model: the bypass is sound because $K \simeq U$
(bisimulation), and it is beneficial because of the boundary-count
inequality. Neither claim is an engineering opinion.
\end{remark}

\section{Cache soundness as bounded staleness}\label{sec:cache}

\setcatalog{61}
\begin{theorem}[Cache soundness]\label{thm:61}\label{thm:cache}
Let $M$ be a molecule whose behavior depends on a state $S$ that
changes only at \emph{event} times (file writes, configuration
reloads, rule swaps), and let $C$ be the caching realization of $M$
that serves the normal form of the state at time $t$ for $\Delta$ time
units. Then $C$ is behaviorally equivalent to $M$ on every request in
$(t, t+\Delta]$ if and only if no event occurs in $(t, t+\Delta]$.
Consequently a sound TTL is any $\Delta$ not exceeding the time to the
next event; the TTL of a cache is a theorem about the event process,
not a tuning constant.
\end{theorem}

\begin{proof}
The cached molecule and the live molecule agree on inputs exactly when
their state agrees, because the output is a deterministic function of
the state and the input (Mealy determinism, \autoref{ch:mol}). The
state agrees on $(t, t+\Delta]$ exactly when no event changed it in
that window; the equivalence is then bisimulation on the window
domain, and the ``if and only if'' is the two directions of that
agreement.
\end{proof}

\begin{remark}
The time-to-live values of the static module in \autoref{ch:situations}
were chosen by intuition; \autoref{thm:cache} makes the choice a
theorem: the minimum time to the next event, or the empirical
event-gap quantile. The invalidation atoms of the control surface
(configuration reload, cache invalidate) are exactly the event process
of the theorem.
\end{remark}

\section{The scheduler estimate is a band contraction}\label{sec:scheduler}

The integer admission scheduler of a production HTTP server maintains
a demand estimate $d' = d + n - d/8$ over $\mathbb{N}$ (a right shift
by 3 in the shipped code). This is the state update of the
fixed-iteration loop of \autoref{ch:trace}; the following theorem makes
its tuning a proof.

\setcatalog{62}
\begin{theorem}[Band contraction]\label{thm:62}\label{thm:ema}
Let $F (x) = x + n - x/8$ on $\mathbb{N}$ with fixed demand
$n \geq 1$. Then:
(1) the band $B = [8n, 8n+7]$ is the fixed set: $F (x) = x$ for
$x \in B$, and every fixed point lies in $B$;
(2) below the band $F$ strictly increases and stays at or below $8n$;
above the band $F$ strictly decreases;
(3) $F$ is nondecreasing, so the iterate $F^{k} (x_0)$ is monotone in
$x_0$;
(4) consequently the iterate converges to $B$ in finitely many steps,
and the crossing time is bounded by the iteration bound of
\autoref{ch:trace} with rate $7/8$: the estimate is a $7/8$
contraction in the continuous relaxation, and the integer rounding
leaves a band of width $8$ instead of a point.
\end{theorem}

\begin{proof}
(1): for $8n \leq x \leq 8n+7$, integer division gives $x/8 = n$, so
$F (x) = x + n - n = x$. Conversely, $F (x) = x$ forces $n = x/8$,
which by the meaning of integer division means $8n \leq x \leq 8n+7$,
that is, $x \in B$.
(2): write $x = 8q + r$ with $q = \lfloor x/8 \rfloor$ and
$0 \leq r \leq 7$, so $F (x) = 7q + r + n$. If $x < 8n$ then
$q \leq n-1$, hence $F (x) \leq 7(n-1) + 7 + n = 8n$, with equality at
$x = 8n-1$. Also $q \leq n-1$ gives $\lfloor x/8 \rfloor \leq n-1$, so
$F (x) = x + n - \lfloor x/8 \rfloor \geq x + 1 > x$. Thus below the
band $F$ strictly increases and $F (x) \leq 8n$. If $x > 8n+7$ then
$q \geq n+1$, so $F (x) = x + n - q \leq x - 1$, and
$F (x) \geq 8q + n - q = 7q + n \geq 7(n+1) + n = 8n+7$. Thus above
the band $F$ strictly decreases and stays at or above $8n+7$. Entry is
finite because each step moves by at least $1$ on $\mathbb{N}$.
(3): if $x \leq y$ then $x/8 \leq y/8$ (integer division is monotone),
so $F (x) = x + n - x/8 \leq y + n - y/8 = F (y)$.
(4): by (2), the distance to the band strictly decreases at each step
outside it, and the state is integer, so the crossing is finite. For
the continuous relaxation $y' = \alpha y + n$ with $\alpha = 7/8$, the
fixed point is $y^{*} = n/(1-\alpha) = 8n$, and
$y' - 8n = \alpha (y - 8n)$, an exact contraction with rate
$\alpha = 7/8$; the iteration bound of \autoref{ch:trace} applies with
that rate, and the integer rounding leaves the band $[8n, 8n+7]$
instead of the point $8n$.
\end{proof}

\begin{remark}
The firewall thresholds and the exponential-decay roll-off of the
same scheduler have the same contractive shape with rates $7/8$ and
$15/16$. \autoref{thm:ema} gives the constants a theorem: the estimate
is within $\varepsilon$ of its attractor after the iteration bound of
\autoref{ch:trace} steps, and the fast-path shortcut of the admission
gate, checking the global bound before the per-source table, is a
cost-monotone rewrite: cheaper, and behavior-preserving on the
rejected set.
\end{remark}

\section{Zero-boundary loops: the hot-path effect discipline}\label{sec:zero-boundary}

\setcatalog{63}
\begin{theorem}[Zero-boundary loops]\label{thm:63}\label{thm:zero-boundary}
Let $L_k (M)$ be the fixed-iteration loop whose body $M$ is a molecule
typed in the affine system of \autoref{ch:linear} with no effectful
atoms: no syscall atom and no allocation atom in its signature. Then
in steady state $\beta (L_k (M)) = 0$: the loop crosses the kernel
boundary only in its setup, and its steady-state cost is purely
userspace, $c (L_k (M)) = k \cdot c (M)$ by the additivity of the cost
law.
\end{theorem}

\begin{proof}
The linear typing of \autoref{ch:linear} forces exact-once consumption and no allocation atom; a
signature without syscall atoms performs no syscalls by the definition
of the boundary count; the loop is fixed-iteration, so no boundary
crossing occurs between iterations; the cost identity is the additive
law applied $k$ times.
\end{proof}

\begin{remark}[The incident as a type error]
The syscall-dominance measured throughout this thesis and the $5.2$
multiplexing gain of the boundary remark are the same fact seen twice:
a hot path that crosses the kernel boundary pays for it per request.
\autoref{thm:zero-boundary} is the enforcement: a dataplane module
whose loop body is not typed zero-boundary is rejected, and the
process-status governor incident is a typing error, not a performance
bug. This is the linearity discipline of \autoref{ch:linear} extended
from allocation to syscalls.
\end{remark}

\section{Route matching as a deterministic automaton}\label{sec:route-automaton}

A production HTTP server's rules are a finite set of \emph{disjoint}
patterns: exact paths and trailing-prefix patterns, each with a method
mask and an exclude list. Disjointness is a load-time theorem (an
overlap check rejects conflicting rule sets before any request is
served). The data is a regular language, and the obvious matcher, a
linear scan over the rules, is only one of the two canonical
realizations.

\setcatalog{64}
\begin{theorem}[Route languages are regular; the pattern trie is a deterministic automaton]\label{thm:64}\label{thm:route-dfa}
Let $P$ be a set of disjoint patterns over the alphabet of path bytes,
each pattern either exact or a trailing prefix. Define the
\emph{pattern automaton} $T (P)$ as the trie whose nodes are the
distinct pattern prefixes; each exact pattern terminates at its own
bytes, each prefix pattern terminates at the prefix separator; every
node carries its include terminals (rule, method mask) and exclude
terminals (rule), split by whether the terminal is exact or prefix.
Then $T (P)$ decides membership of a path $p$ in the route language of
$P$ in exactly $|p|$ transitions, without backtracking and without
allocating.
\end{theorem}

\begin{proof}
Walk $p$ byte by byte. A prefix include terminal fires the moment its
node is visited, because a prefix pattern matches $p$ if and only if
$p$ starts with the prefix, which is exactly the walk reaching that
node. An exact terminal fires only at the last visited node, because
an exact pattern matches $p$ if and only if the walk ends there.
Excludes are collected into a running bitmask: prefix excludes set
their bit when visited, exact excludes only at the final node.
Disjointness guarantees at most one rule can match $p$, so the first
method-matching include terminal is the candidate, and the candidate
matches if and only if no exclude bit of it was set. Both firing rules
are exactly the predicate of the linear matcher, so the automaton and
the scan agree on every path. The loop is one pass over the bytes:
$|p|$ transitions, each a binary search over the node's sorted
children. The equivalence is machine-checked in-crate by a
200{,}000-path fuzz test comparing the automaton against the scan on
every path, including prefix boundaries, method splits, and exclude
vetoes.
\end{proof}

\setcatalog{65}
\begin{theorem}[The automaton is the Myhill--Nerode refinement]\label{thm:65}\label{thm:route-mn}
The states of $T (P)$ are the distinct pattern prefixes. Under the
Myhill--Nerode quotient of the route language, two prefixes are
equivalent exactly when their futures, the sets of patterns that can
still complete from them, coincide; the trie states are the prefixes
that carry distinct futures. Hence $T (P)$ is a refinement of the
minimal deterministic automaton of the route language, and it is
exactly minimal for prefix-closed pattern sets with pairwise-distinct
futures.
\end{theorem}

\begin{proof}
The Myhill--Nerode relation identifies prefixes with identical right
languages \citep{nerode}; the quotient is the minimal deterministic
automaton \citep{hopcroft}, and Brzozowski's derivatives compute it
algebraically \citep{brzozowski}. For a trie of patterns, the future
of a node is determined by the set of terminals in its subtree: two
nodes with identical terminal sets are identified by the quotient, and
nodes with different terminal sets are distinguished by the suffix
that ends at a terminal present in one but not the other. The trie
therefore collapses exactly the same states as the quotient when
futures are distinct, and refines it otherwise. Prefix languages have
minimal automata of exactly this shape: one state per distinct prefix
\citep{arxiv-syntactic}. For the shipped rule set the futures are in
practice always distinct, so the automaton is the minimal one.
\end{proof}

\section{Realization choice is a measured threshold}\label{sec:route-crossover}

\setcatalog{66}
\begin{theorem}[Hybrid matcher selection]\label{thm:66}\label{thm:hybrid}
Let the scan cost be affine in the rule count,
$c_{scan} (R) = aR + b$, and the automaton cost be flat in $R$. Then by
the affine crossover theorem (\autoref{thm:crossover}) the cheaper
realization is decided by a single threshold $R^{\ast}$, and on the
measured grid the threshold is the lattice minimum of
\autoref{thm:lattice}.
\end{theorem}

\begin{proof}
This is \autoref{thm:crossover} instantiated with payload $R$. The
measured table (release build, nanoseconds per match on identical rule
sets and paths):
\[
\begin{array}{l|rrrr}
R & 2 & 4 & 8 & 16\\\hline
\text{scan} & 5.9 & 10.6 & 25.8 & 59.6\\
\text{automaton} & 29.4 & 31.7 & 31.8 & 43.9
\end{array}
\]
The scan is affine in $R$ (fit $\approx 4.2R - 2.5$\,ns); the
automaton is flat. The fitted crossover is $R^{\ast} \approx 10$--$12$,
and on the measured grid $\{2,4,8,16,\ldots\}$ the automaton first
wins at $R = 16$. The shipped constant is the lattice minimum of
\autoref{thm:lattice}: the decision rule, scan below 16 rules and
automaton at 16 or above, is correct because the scan is cheaper for
every $R \leq 8$ (the production range) and the automaton is cheaper
for every $R \geq 16$.
\end{proof}

\begin{remark}
The same threshold method selects between two realizations of the
reactor loop at startup (\autoref{sec:autotune}) and between the byte
path and the file-transfer path of a static-file molecule
(\autoref{sec:threshold}). In each case the choice is computed from
measured cost lines, not chosen by taste.
\end{remark}

\section{The zero-allocation claim is machine-checked}\label{sec:alloc-count}

The hot-path discipline of \autoref{thm:zero-boundary} is a typing
claim. The claim is enforced by measurement: a counting allocator,
installed as the global allocator of a test process, counts every
allocation the process makes. The datapath test runs a full echo
cycle (receive batch, validate, echo batch) for 50 rounds and asserts
that the counter is unchanged. The control test, which allocates a
known buffer, proves the counter observes allocations; without the
control the assertion could pass vacuously.

\section{Startup lattice autotuning}\label{sec:autotune}

The reactor strategy is a realization choice with the same shape as
\autoref{thm:hybrid}: each strategy has a fixed cost and a
load-dependent cost, and the cheapest strategy depends on the machine
and the kernel. The startup autotuner measures the candidates on the
running machine with a fixed benchmark (windowed echo), applies a
soundness floor (a candidate whose TCP path cannot sustain half the
best TCP rate is disqualified, because a reactor that stalls its
connection path is not a viable production choice), and selects the
remaining candidate with the highest UDP throughput. On the reference
machine the measured order is: readiness-driven first, completion-driven
second, completion-driven with a kernel polling thread last; the
completion-driven candidates are disqualified on this kernel by the TCP
soundness floor. The winner is a lattice minimum: the finite set of
strategy realizations, graded by the measured boundary-enriched cost,
with the soundness floor as the admissibility predicate.

\section{Measured case studies: before and after}\label{sec:demos}

This section applies the discovery method of \autoref{sec:lattice} to
two concrete molecules. For each molecule the section defines a
baseline realization, applies individual optimizations one at a time,
and finally applies the full stack. Every number is measured on one
machine (Intel Core i5-5200U, kernel 7.2.0\_1) with identical inputs
across stages; the raw output is in the source distribution
(\texttt{Docs/paper/data/demo-metrics.csv}). The metric set is:
instructions and instructions per cycle (IPC), branch miss rate, L1
data-cache miss rate, L3 (last-level cache, LLC) load miss rate, dTLB
miss rate, wall time (latency, with p50 and p99 percentiles over 20
runs per stage), peak resident set (RSS), and page faults. The
processor does not expose the stalled-cycles performance events, so
pipeline stalls are reported through the cycle and IPC columns, which
incorporate them.

\subsection{Case study 1: the input-output molecule}

The molecule copies a 512\,MiB file and computes a byte-sum checksum
over the data, a miniature of a file-serving pipeline. The baseline
realization is the naive atom wiring.

\begin{algorithm}[t]
\caption{Baseline: the input-output molecule (before).}
\begin{algorithmic}[1]
\For{each block}
  \State $buf \gets \Call{alloc}{block\_size}$ \Comment{per-block allocation}
  \State $n \gets \Call{read}{src, buf}$ \Comment{one syscall}
  \State $\Call{write}{dst, buf[0..n]}$ \Comment{one syscall}
  \State $acc \gets acc + \Call{checksumScalar}{buf[0..n]}$ \Comment{byte loop}
\EndFor
\end{algorithmic}
\end{algorithm}

The first optimization fuses the boundary: a preallocated buffer and
gathered system calls remove the per-block allocation and the two
syscalls per block become one system call per transfer.

\begin{algorithm}[t]
\caption{Optimization 1: preallocation and gathering (after).}
\begin{algorithmic}[1]
\State $buf \gets \Call{alloc}{block\_size}$ \Comment{once}
\While{$n \gets \Call{readv}{src, buf} > 0$}
  \State $\Call{writev}{dst, buf[0..n]}$
  \State $acc \gets acc + \Call{checksumScalar}{buf[0..n]}$
\EndWhile
\end{algorithmic}
\end{algorithm}

The second optimization fuses the compute: the checksum runs in the
same pass over the same buffer, vectorized over 32-byte blocks with a
scalar remainder, so the data is touched once and the instruction
count drops.

\begin{algorithm}[t]
\caption{Optimization 2: fused vectorized checksum (after).}
\begin{algorithmic}[1]
\State $buf \gets \Call{alloc}{block\_size}$
\While{$n \gets \Call{readv}{src, buf} > 0$}
  \State $\Call{writev}{dst, buf[0..n]}$
  \State $acc \gets acc + \Call{checksumSimd}{buf[0..n]}$ \Comment{same pass}
\EndWhile
\end{algorithmic}
\end{algorithm}

\begin{table}[t]
\centering
\caption{Input-output molecule: measured stages. Baseline (B), plus
gathering ($B+G$), plus fused vectorized checksum ($B+G+V$).}
\label{tab:demo1}
\begin{tabular}{lrrr}
\toprule
Metric & B & $B+G$ & $B+G+V$\\
\midrule
Throughput (MiB/s) & 1381.2 & 1405.1 & 1593.9\\
Wall time, p50 (ms) & 357 & 352 & 307\\
Wall time, p99 (ms) & 447 & 442 & 411\\
Instructions & 2.179 G & 2.191 G & 2.062 G\\
Cycles & 1.917 G & 1.870 G & 1.821 G\\
IPC & 1.14 & 1.17 & 1.13\\
Branch misses (\%) & 0.40 & 0.40 & 0.41\\
L1 data misses (\%) & 11.66 & 10.19 & 11.01\\
L3 load misses (\%) & 34.9 & 32.5 & 29.1\\
dTLB misses (\%) & 0.011 & 0.010 & 0.011\\
Page faults & 90 & 91 & 91\\
Peak RSS (KiB) & 2016 & 2076 & 2028\\
\bottomrule
\end{tabular}
\end{table}

The stacking is compounding: gathering alone gains 1.7\%; the fused
vectorized checksum on top of gathering gains another 13.4\%, for a
cumulative 15.4\% over the baseline. Wall time falls 14.0\% at the p50
(357 to 307 ms) and 8.1\% at the p99 (447 to 411 ms). The instruction
count falls 5.4\% and the cycle count falls 5.0\%; the throughput gain
exceeds the cycle gain because the vectorized checksum reduces memory
traffic (L1-dcache loads fall from 647 M to 611 M, and the L3 load
miss rate falls from 34.9\% to 29.1\%).

\subsection{Case study 2: the pure compute molecule}

The molecule enumerates the primitive Pythagorean triples
$(a, b, c)$ with $c \leq N$: triples with $a^2 + b^2 = c^2$ and
$\gcd (a, b) = 1$. The baseline is the naive wiring: a double loop
over $a$ and $b$ with a square-root test per candidate and a
per-candidate allocation.

\begin{algorithm}[t]
\caption{Baseline: the triple generator (before).}
\begin{algorithmic}[1]
\For{$a \gets 1$ \textbf{to} $N$}
  \For{$b \gets a$ \textbf{to} $N$}
    \State $c \gets \lfloor \sqrt{a^2+b^2} \rfloor$ \Comment{float per candidate}
    \If{$c^2 = a^2+b^2$ \textbf{and} $c \leq N$ \textbf{and} $\gcd (a,b) = 1$}
      \State $\Call{emit}{a, b, c}$ \Comment{per-triple allocation}
    \EndIf
  \EndFor
\EndFor
\end{algorithmic}
\end{algorithm}

Optimization 1 replaces the search by Euclid's formula: every
primitive triple arises exactly once from coprime parameters $m > n$
with $m - n$ odd, as $(m^2 - n^2, 2mn, m^2 + n^2)$. No square root
and no allocation remain.

\begin{algorithm}[t]
\caption{Optimization 1: Euclid's formula (after).}
\begin{algorithmic}[1]
\For{$m \gets 2$ \textbf{while} $m^2 + 1 \leq N$}
  \For{$n \gets 1$ \textbf{to} $m - 1$}
    \If{$(m-n)$ odd \textbf{and} $c = m^2+n^2 \leq N$ \textbf{and} $\gcd (m,n) = 1$}
      \State $\Call{emit}{(m^2-n^2,\; 2mn,\; c)}$
    \EndIf
  \EndFor
\EndFor
\end{algorithmic}
\end{algorithm}

Optimization 2 unrolls the inner sweep: two adjacent candidates are
evaluated per iteration, which halves the loop control overhead for
the dominant parity class.

\begin{table}[t]
\centering
\caption{Triple generator: measured stages. Baseline (B), Euclid
(E), Euclid unrolled ($E+U$).}
\label{tab:demo2}
\begin{tabular}{lrrr}
\toprule
Metric & B & E & $E+U$\\
\midrule
Input & $N = 20{,}000$ & $N = 2{,}000{,}000$ & $N = 2{,}000{,}000$\\
Triples found & 3{,}186 & 318{,}320 & 318{,}320\\
Instructions per triple & 2.04 M & 2.29 k & 2.01 k\\
Instructions (total) & 6.483 G & 0.728 G & 0.640 G\\
Cycles (total) & 3.333 G & 1.214 G & 1.159 G\\
IPC & 1.95 & 0.60 & 0.55\\
Wall time, p50 (ms) & 1215 & 23 & 22\\
Wall time, p99 (ms) & 1223 & 23 & 23\\
Branch misses (\%) & 0.02 & 5.40 & 6.00\\
L1 data misses (\%) & 7.42 & 6.92 & 6.77\\
L3 load misses (\%) & 43.4 & 35.5 & 54.0\\
dTLB misses (\%) & 0.21 & 0.36 & 0.28\\
Page faults & 78 & 77 & 77\\
Peak RSS (KiB) & 1936 & 1916 & 1916\\
\bottomrule
\end{tabular}
\end{table}

The comparison is normalized per triple: Euclid's formula uses 2.29
thousand instructions per triple against 2.04 million for the brute
force, a factor of about 890. The wall time per triple (p50 over 20
runs) falls from about 381\,$\mu$s to about 72\,ns, a factor of about
5{,}300 (the brute force tests every pair up to $N$; Euclid's formula
tests only the coprime parameter pairs). The unroll on top of Euclid
cuts the instruction count a further 12\% and the cycle count a
further 5\%, and the wall time at the p50 falls from 23 to 22 ms. The
branch miss rate rises with Euclid because the parity and coprimality
tests are genuinely data-dependent branches; the absolute instruction
count is what dominates.

\subsection{How the optimizations stack}

For the input-output molecule the stack is
baseline $\to$ gathering $\to$ fused vectorized checksum, with
cumulative improvement 1.7\% then 15.4\%. For the triple generator the
stack is baseline $\to$ Euclid $\to$ unroll, with per-triple
instruction reduction of factor 890 then a further 12\%. In both cases
the individual gains compound: each stage preserves the previous
stage's behavior (the checksum and the triple set are identical across
stages) and adds a cost reduction on top. This is the lattice method
of \autoref{sec:lattice} in measurement form: each stage is a
behavior-preserving rewrite, and the stages compose.

\section{Verification and summary}\label{sec:breakthrough-verify}

Every arithmetic claim of this chapter is proved in full in the text:
the boundary laws (\autoref{thm:boundary}), the batch identities
(\autoref{thm:batch}), the affine crossover (\autoref{thm:crossover}),
the lattice finiteness (\autoref{thm:lattice}), the band contraction
(\autoref{thm:ema}), and the automaton equivalence
(\autoref{thm:route-dfa}) and its Myhill--Nerode refinement
(\autoref{thm:route-mn}). The computational claims, the finite
checks that the arithmetic licenses, are discharged by the executable
verifiers of \autoref{app:a}; the measured tables of this chapter are
reproducible from the source distribution.

The method of this chapter: enrich the cost model with the kernel
boundary, prove that realization choice is a threshold, and take the
minimum of the finite lattice of behavior-preserving rewrites. It is
the answer to the question ``what is the last optimization?'' The
method is a decision procedure that
\autoref{thm:lattice} proves exists, applied to the molecules this
thesis constructs. The optimizations discovered this way are exactly
the ones the measurements confirm: fewer syscalls per request, the
correct datapath at the correct payload size, and a zero-boundary loop
as the type-correct realization of a hot path.
